\chapter{Diseño de la API REST (Go)}

El servidor en \go{cmd/simulation/main.go} expone los endpoints sobre \go{http.ListenAndServe(":8080", mux)}.

\begin{table}[H]
  \centering
  \caption{Endpoints de la API REST}
  \begin{tabular}{@{}llp{6.5cm}@{}}
    \toprule
    \textbf{Método} & \textbf{Ruta} & \textbf{Descripción} \\
    \midrule
    GET & \texttt{/api/tree} & Estado del árbol, métricas e \texttt{inOrder} \\
    POST & \texttt{/api/products} & Insertar o actualizar producto \\
    GET & \texttt{/api/products/\{id\}} & Búsqueda por ID en el árbol \\
    DELETE & \texttt{/api/products/\{id\}} & Eliminar producto \\
    POST & \texttt{/api/reset} & Reiniciar dataset semilla \\
    \bottomrule
  \end{tabular}
\end{table}

\section{Ejemplo de respuesta con traza}

\begin{lstlisting}[style=jsoncode,caption={Fragmento de respuesta POST /api/products}]
{
  "message": "Producto 109 insertado",
  "trace": {
    "path": [105, 109],
    "inserted": true,
    "depth": 1,
    "maxDepth": 3,
    "scapegoat": null,
    "rebuiltKeys": []
  }
}
\end{lstlisting}

\section{Sincronización y concurrencia}

\begin{lstlisting}[style=gocode,caption={Estado compartido protegido con mutex}]
type appState struct {
    mu   sync.Mutex
    tree *scapegoat.Tree[int, db.Product]
    sql  *sql.DB
    mode string
}
\end{lstlisting}

Cada handler adquiere \texttt{mu} antes de mutar o leer \texttt{tree}. Las operaciones SQL usan \texttt{context.WithTimeout} de 10 segundos.

\begin{figure}[H]
  \centering
  \begin{tikzpicture}[
    actor/.style={rectangle, draw, rounded corners, minimum width=1.8cm, align=center, font=\footnotesize},
    arrow/.style={-{Stealth[length=2mm]}, thick}
  ]
    \node[actor] (u) {Usuario};
    \node[actor, right=1cm of u] (v) {Vue};
    \node[actor, right=1cm of v] (a) {API};
    \node[actor, right=1cm of a] (t) {Árbol};
    \draw[arrow] (u) -- (v);
    \draw[arrow] (v) -- node[above, font=\tiny]{POST} (a);
    \draw[arrow] (a) -- node[above, font=\tiny]{InsertWithTrace} (t);
    \draw[arrow] (t) -- node[below, font=\tiny]{trace JSON} (v);
  \end{tikzpicture}
  \caption{Secuencia de inserción vía API}
  \label{fig:seq-api}
\end{figure}
